\Needspace{5\baselineskip}
\paragraph{Illustrative annotations.}
Figure~\ref{fig:sktl-sample} illustrates complete boundaries and contextual typing. A further handbook case, 熟悉Linux基本操作和原理, selects Linux基本操作 (S) and 原理 (K): the second span refers back to Linux without copying it into the annotation. This Chinese coordination pattern requires separate boundary and type decisions (handbook R14, case 130).
\begin{figure}[htbp]
\centering\includegraphics[width=\linewidth]{figures/boundary_context_round2.pdf}
\caption{Examples of boundary completeness, coordination, and contextual type. Highlighted substrings are selected spans; surrounding predicates remain available as context. These teaching examples illustrate the current handbook and do not revise historical annotations. Type codes and English glosses are display aids.}
\label{fig:sktl-sample}
\end{figure}
